\section{Dara graph applications and the authority to change a causal model}
\label{sec:exec25-dara-reuse}
Dara provides a concrete application and graph-editing reference for the
proposed causal library. The inspected repository revision
\texttt{0ca24611ef35a9725f14a}\allowbreak\texttt{f6cda83718482665294}
is dated 30 September 2026. Its selected Python core, component and
application-generator packages, and JavaScript graph editor, declare version
1.29.11. Nineteen retained source files were authenticated against the complete
Git tree. The root and editor notices refer to Apache 2.0; dependencies,
assets, deployed services and commercial offerings remain separate rights and
qualification questions. This is an inspection of immutable source bytes,
not an installation or demonstration of a causal service.
\citep{exec25_dara_source}

The causaLens/Dara umbrella is resolved into distinct proposal decisions.
Dara's framework, graph viewer/editor and application CLI are retained as
references, with component-specific conditional integration. Chapter~\ref{scisyn:software} separately identifies the later advertised
causaLens service offering. Its commercial integration remains deferred until
a supported artifact, interface, applicable terms and substantive technical
evidence are selected; the open Dara source establishes none of those by itself.
Open framework code cannot establish proprietary product capabilities,
entitlement, pricing or operational acceptance. The proposal therefore neither
borrows commercial approval from the framework nor infers a commercial
absence from the selected open-source paths.~\citep{exec25_dara_choice}

A branch manager may want to revise a diagram after learning that an account
status was recorded after a servicing call. A useful interface can display the
graph and collect a proposed correction. Scientific acceptance additionally
requires the question version, source evidence, variable identity, observation
clock, edit authority and consequences for identification. A drag action or
confirmation flag is evidence of a user interaction; it does not prove that
the revised arrow represents a causal mechanism. This distinction is central
to using an application framework around the proposed causal engine.
\citep{exec25_dara_choice}

\subsection{Graph modes and conversions are scientific contracts}
The selected Python viewer accepts a causal-graph object, skeleton, dictionary
or variable containing graph data. Its local validator returns variables
unchanged and checks the declared input type for other inputs. The same class
changes a non-default editor mode to DEFAULT when a planar layout is selected,
with a warning. The source describes DEFAULT as a DAG-oriented viewer, PAG
mode as displaying additional endpoint types, and RESOLVER as supporting
edge acceptance. A proposed adapter must make such mode changes visible and
verify that the declared graph class and unresolved marks survive conversion.
A type check on a dictionary is not a full proof of causal-graph validity or
of the assumptions behind the displayed graph.~\citep{exec25_dara_source}

The inspected editor types distinguish directed, bidirected, undirected and
partially marked edges, including an internal backwards-directed form. These
symbols can carry materially different scientific meanings. A latent-variable
analysis that returns unresolved endpoint marks must not be displayed as one
fully directed intervention graph simply because a layout or editing mode
expects it. A round-trip oracle should compare identifiers, endpoint marks,
lag or time information where admitted, graph class, metadata provenance and
the target query, not just whether the picture contains the same number of
edges. Unsupported graph classes should produce an explicit refusal or a
clearly labeled visualization-only view.~\citep{exec25_dara_source,exec25_dara_choice}

One selected serialization composition illustrates the need for that oracle.
The edge serializer converts a backwards-directed representation to a directed
edge and exchanges its embedded source and destination. The enclosing graph
serializer again exchanges those embedded endpoints while placing the edge
under reversed dictionary keys. For a stipulated internal pair $(A,B)$, the
first exchange yields embedded $(B,A)$; the second yields $(A,B)$, while the
outer keys are $(B,A)$. This is an exact symbolic trace of the selected source
branches, not a claim that an installed frontend has produced a reachable
failure. It identifies a concrete consistency test: the endpoint identifiers
inside every serialized edge must agree with its enclosing source/target
keys under the declared convention. Both direct and backwards representations
must pass the same canonicalization test.~\citep{exec25_dara_source,exec25_dara_choice}

\subsection{Editing, acyclicity and evidence authority}
The default-mode edge-add path checks whether the proposed addition creates
a cycle. In the selected reverse-edge path, the original edge is removed from
a graph copy and the original selected tuple is supplied to the cycle helper;
the eventual reducer exchanges the edge direction. That argument sequence
requires a targeted test against the intended reversed graph. An authored
example starts from $A\to B$, $A\to C$, $C\to B$, which is acyclic. Reversing
$A\to B$ creates $B\to A\to C\to B$. Checking the original direction does
not certify the reversed one. The source trace and exact graph calculation
motivate a qualification fixture; no upstream code or native failure was
executed. The test must also respect mode-specific semantics: DAG acyclicity
is not an all-purpose validity criterion for every partial ancestral graph.
\citep{exec25_dara_source,exec25_dara_choice}

The selected reducer records edge acceptance in rendering metadata. The
Python/JavaScript wrapper updates its graph variable and invokes the update
action, while the editor coalesces outgoing updates with a trailing
150-millisecond debounce. These mechanisms support interaction, but a
coalesced callback does not by itself provide an authenticated record of every
scientific edit or a transaction approved by a causal backend. The proposed
integration should distinguish a local draft, a submitted change and an
accepted model version. A change prepared for graph version $g_1$ must be
reconciled or refused if the current graph is $g_2$; it must not silently
apply to a different model.~\citep{exec25_dara_source,exec25_dara_choice}

For the servicing example, an edit that removes a mediator from a proposed
adjustment set changes the analysis contract even if the screen still shows
the same nodes. The accepted record should identify the author or authorized
role, evidence, old and new graph digests, question version, restriction type,
reason and resulting identification status. An LLM can draft the explanation
or propose alternatives, while the service preserves unknown versus absent
edges and conditional versus identified effects. Rendering acceptance,
authority to edit and scientific justification are three separate fields.
\citep{exec25_dara_choice}

\subsection{Application CLI versus a causal service interface}
The retained metadata declares application-creation and application-runtime
CLI entry points. Their role is useful for packaging an interface; it does
not establish a CLI for causal estimands or an MCP service contract. The
selected component dependency on \texttt{cai-causal-graph>=0.3.6} also leaves
a concrete delegate and environment to qualify. Python and JavaScript
package versions, bundled assets, protocol schemas and dependencies must be
bound together before a deployed interface can be called reproducible.
A root notice does not settle the entire dependency or deployment inventory.
\citep{exec25_dara_source}

Conditional integration would place a versioned graph editor in front of the
library's question, identification and execution contracts. Requests would
carry graph class, identifiers and marks, question version, data/schema
digests, restrictions, edit provenance, authority, expected graph version and
a deadline. Responses would separate draft rendering from accepted graph
state, identified or conditional effects, refusal and execution status. A
port of selected editor code would still require rights, dependency and
semantic tests; it would not inherit a native causal guarantee from the
original UI. Direct adoption as a causal inference engine is not supported
by the selected application source.~\citep{exec25_dara_choice}

The proposed public qualification suite uses finite known graphs and generated
structural laws for scientific truth. It includes all admitted edge marks,
identifier permutations, backwards-edge serialization, the reversal-cycle
example, an attempted planar/PAG mode change, stale edits, rejected backend
changes, duplicate updates and late callbacks. A fixed graph-file workflow,
a permitted frozen Dara interface and the proposed adapter should receive
the same questions and information. The protocol must freeze schemas,
versions, graph conventions, scientific acceptance and complete denominators
before native measurement. Permitted public data can add realistic schema
and interaction stress without supplying unobserved causal truth. None of
these native trials is claimed as completed here.~\citep{exec25_dara_choice}

Scientific owners would approve graph restrictions, questions and
identification; interface owners would preserve rendering and serialization;
service owners would enforce version and authority checks; rights/platform
owners would select permitted artifacts; evaluators and maintainers would
control fixtures and regressions. These remain proposed responsibilities.
An authored USD 1,000 preparation cost, USD 1,000 of all attempts,
USD 2,500 scientific review and USD 500 maintenance allocation sum to
USD 5,000. For 100 attempts, 80 completions and 50 admissible outputs, costs
are USD 50, USD 62.50 and USD 100. These are accounting examples, not vendor
prices. The selected reference route remains conditional on rights,
semantic preservation and measured useful scientific outcomes.
\citep{exec25_dara_choice}
